For odd $n\ge3$, a putative square $m^2$ has odd $m$. Then $(\dfrac{m-n}{2}) (\dfrac{m+n}{2}) = 2^{n-2}$, and the two factors differ by odd $n$. The odd factor must be 1, necessarily the smaller, giving $2^{n-2}=n+1$. This fails at $n=3$, and for odd $n\ge5$ the left side exceeds the right (induct by increments of two). The case $n=1$ fails directly. For even $n\ge14$, $n^2 < 2^{\dfrac{n}{2} + 1} + 1$, so $2^n+n^2$ lies strictly between $(2^{\dfrac{n}{2}})^2$ and $(2^{\dfrac{n}{2}} + 1)^2$. The inequality holds at 14 and persists after adding two, since the exponential doubles while $(\dfrac{n+2}{n})^2 < 2$. Checking $n=2,4,6,8,10,12$ gives $8,32,100,320,1124,4240$; only 100 is square. Therefore $\boxed{n=6}$.
